% ID: SEL-D0061-OA-E01
% ESTADO: validada. S061.
\textbf{a)} Para cada $y\in[0,6]$, la región va desde $x=0$ hasta la menor de las cotas $2y$ y $15-y$. Ambas fronteras se incluyen.
\begin{center}
\setlength{\unitlength}{0.9mm}
\begin{picture}(96,59)(-10,-8)
\qbezier(0,0)(39,0)(78,0)
\qbezier(0,0)(0,23)(0,46)
\qbezier(0,0)(35,17.5)(70,35)
\qbezier(70,35)(66.5,38.5)(63,42)
\qbezier(63,42)(31.5,42)(0,42)
\put(80,0){$x$}\put(2,46){$y$}
\put(71,33){\scriptsize $(10,5)$}
\put(59,44){\scriptsize $(9,6)$}
\put(-10,40){\scriptsize $(0,6)$}
\put(-7,-6){\scriptsize $(0,0)$}
\end{picture}
\end{center}
\textbf{b)} En $y=0$ sólo cabe $x=0$. En $y=6$ quedan los extremos $(0,6)$ y $(9,6)$. Las cotas oblicuas se igualan cuando $2y=15-y$, lo que da el cuarto vértice $(10,5)$.

\textbf{c)} En $(0,0),(0,6),(9,6),(10,5)$ los valores de $F$ son $0,30,102,105$. Como la función es lineal y sólo un vértice alcanza el mayor valor, el máximo es $\boxed{105}$ únicamente en $(10,5)$.
